% practical.tex
\section{Practical Considerations}
\label{sec:practical}

We discuss practical aspects of deploying KnowFeat based on our experience with the SimECNY dataset.

\textbf{Computational Cost.}
The full pipeline takes $\sim$3.2 hours for 4 rounds on SimECNY (52 candidates, 18 accepted). L3 model evaluation dominates at 41\% of total time ($\sim$2.5 min per candidate), followed by LLM reasoning at 23\% ($\sim$52 sec per candidate). Progressive filtering saves $\sim$53 minutes by rejecting 21 candidates at L1/L2 before expensive L3 evaluation. The LLM API cost is approximately \$2.80 for the full run (DeepSeek-V4-0324-Flash). Importantly, this is a one-time setup: at inference, only the generated Python code runs, adding $<$50ms per 200K-transaction batch.

\textbf{Compliance Integration.}
Provenance cards map directly to existing AML compliance review processes. The knowledge asset identifiers (e.g., RULE\_006, IND\_15) reference the same regulatory documents that compliance officers work with. Domain experts from the Postal Savings Bank of China reported that the structured format reduced feature review time from 2--3 hours to 15--30 minutes per feature, as the reasoning chain eliminates the need to reconstruct the designer's rationale. This directly supports Article~32 of China's Anti-Money Laundering Law (2024 revision), which mandates documenting the basis for suspicious transaction reporting.

\textbf{Lessons Learned.}
Three insights emerged: (1)~\emph{Knowledge curation is the bottleneck}: compiling our knowledge base (38 indicators, 15 rules, 12 expert opinions, 309 court documents) required $\sim$3 person-weeks; this upfront cost is amortized across all generated features but represents the main barrier for new domains. (2)~\emph{Verification thresholds require domain calibration}: the L3 acceptance criteria (e.g., $\Delta$Recall $> 0$ with Precision $\geq 90\%$) must be set with domain stakeholders. (3)~\emph{Feature maintenance is ongoing}: regulatory requirements evolve, and a periodic refresh cycle (e.g., quarterly) would be necessary in production.
